% Sorgente LaTeX canonico, importato e revisionato dal precedente sorgente Markdown.
\chapter{Limiti di funzione}
\label{chap:7}

\begin{obiettivocapitolo}{scegliere il metodo più rapido per ogni forma di limite}
\prioritaalta. La sostituzione diretta è sempre il primo controllo; le tecniche avanzate si usano solo dopo aver riconosciuto una forma indeterminata.
\end{obiettivocapitolo}
\begin{sceltametodo}{forma del limite}
\begin{longtable}{@{}p{.28\textwidth}p{.31\textwidth}p{.33\textwidth}@{}}
\toprule
\textbf{Forma riconosciuta} & \textbf{Primo tentativo} & \textbf{Passaggio successivo}\\
\midrule\endhead
Rapporto di polinomi & confronto dei gradi & raccogli la potenza dominante\\
$0/0$ algebrico & scomposizione e semplificazione & razionalizzazione o limiti notevoli\\
Differenza di radicali & coniugato & raccogli il termine dominante\\
$0/0$ o $\infty/\infty$ regolare & semplificazioni preliminari & de l'Hôpital se le ipotesi sono verificate\\
Prodotto $0\cdot\infty$ & trasformalo in quoziente & confronto, Taylor o de l'Hôpital\\
Differenza $\infty-\infty$ & denominatore comune, razionalizzazione o raccolta & Taylor se c'è cancellazione\\
Potenza $1^\infty$, $0^0$, $\infty^0$ & logaritmo & calcola $\lim g\log f$ e poi esponenzia\\
Funzioni equivalenti senza cancellazione & equivalenze asintotiche & controlla che compaiano in prodotti/quozienti\\
Cancellazione di termini principali & Taylor di ordine superiore & trova il primo coefficiente non nullo\\
Funzione oscillante limitata per fattore nullo & teorema dei carabinieri & stima il modulo\\
\bottomrule
\end{longtable}
\end{sceltametodo}
\begin{proceduraesame}{limite completo}
\begin{enumerate}
\item Identifica dominio locale e verso di avvicinamento.
\item Sostituisci formalmente e classifica il risultato: determinato, infinito con segno, oppure indeterminato.
\item Scegli la trasformazione minima che espone il termine dominante.
\item Se usi equivalenze o Taylor, specifica il regime e conserva l'ordine sufficiente a superare le cancellazioni.
\item Se usi de l'Hôpital, verifica forma, derivabilità e comportamento del rapporto delle derivate.
\item Concludi indicando anche il segno nei limiti infiniti o laterali.
\end{enumerate}
\end{proceduraesame}
\begin{errorefrequente}{equivalenze in somme e differenze}
Da $f\sim g$ non segue in generale $f-h\sim g-h$. Quando i termini principali si cancellano, serve uno sviluppo di ordine superiore.
\end{errorefrequente}

\section{Limite finito in un punto finito}\label{limite-finito-in-un-punto-finito}

Siano \(f:D\subseteq\mathbb R\to\mathbb R\), \(x_0\in D'\) e \(L\in\mathbb R\).

\[
\lim_{\substack{x\to x_0\\x\in D}}f(x)=L
\]

\[
\Longleftrightarrow
\forall\varepsilon>0,\ \exists\delta>0,\ \forall x\in D:\
0<|x-x_0|<\delta
\Rightarrow
|f(x)-L|<\varepsilon.
\]

In questa definizione \(x\) si avvicina a \(x_0\) senza essere obbligato a coincidere con \(x_0\); il valore \(f(x_0)\) può anche non essere definito. Il limite descrive il comportamento \textbf{vicino} al punto, non necessariamente nel punto.

Negazione:

\[
\lim_{\substack{x\to x_0\\x\in D}}f(x)\ne L
\]

\[
\Longleftrightarrow
\exists\varepsilon_0>0,\ \forall\delta>0,\ \exists x\in D:\
0<|x-x_0|<\delta,
\quad
|f(x)-L|\ge\varepsilon_0.
\]

Approfondimento: \href{https://ingegnerismo.it/matematica/limite-di-funzione/}{limite di funzione}.

\section{Regimi finiti e infiniti}\label{regimi-finiti-e-infiniti}

Nelle formule con limite finito si assume \(L\in\mathbb R\).

Per \(x_0\in D'\):

\[
\lim_{\substack{x\to x_0\\x\in D}}f(x)=+\infty
\]

\[
\Longleftrightarrow
\forall M>0,\ \exists\delta>0,\ \forall x\in D:\
0<|x-x_0|<\delta
\Rightarrow
f(x)>M.
\]

\[
\lim_{\substack{x\to x_0\\x\in D}}f(x)=-\infty
\]

\[
\Longleftrightarrow
\forall M>0,\ \exists\delta>0,\ \forall x\in D:\
0<|x-x_0|<\delta
\Rightarrow
f(x)<-M.
\]

Per \(D\) non superiormente limitato:

\[
\lim_{\substack{x\to+\infty\\x\in D}}f(x)=L
\]

\[
\Longleftrightarrow
\forall\varepsilon>0,\ \exists A\in\mathbb R,\ \forall x\in D:\
x>A\Rightarrow|f(x)-L|<\varepsilon.
\]

\[
\lim_{\substack{x\to+\infty\\x\in D}}f(x)=+\infty
\]

\[
\Longleftrightarrow
\forall M>0,\ \exists A\in\mathbb R,\ \forall x\in D:\
x>A\Rightarrow f(x)>M.
\]

\[
\lim_{\substack{x\to+\infty\\x\in D}}f(x)=-\infty
\]

\[
\Longleftrightarrow
\forall M>0,\ \exists A\in\mathbb R,\ \forall x\in D:\
x>A\Rightarrow f(x)<-M.
\]

Per \(D\) non inferiormente limitato:

\[
\lim_{\substack{x\to-\infty\\x\in D}}f(x)=L
\Longleftrightarrow
\forall\varepsilon>0,\ \exists A\in\mathbb R,\ \forall x\in D:\
x<A\Rightarrow|f(x)-L|<\varepsilon.
\]

\[
\lim_{\substack{x\to-\infty\\x\in D}}f(x)=+\infty
\Longleftrightarrow
\forall M>0,\ \exists A\in\mathbb R,\ \forall x\in D:\
x<A\Rightarrow f(x)>M.
\]

\[
\lim_{\substack{x\to-\infty\\x\in D}}f(x)=-\infty
\Longleftrightarrow
\forall M>0,\ \exists A\in\mathbb R,\ \forall x\in D:\
x<A\Rightarrow f(x)<-M.
\]

Approfondimento: \href{https://ingegnerismo.it/matematica/limite-di-funzione/}{limite di funzione}.

\section{Caratterizzazione sequenziale}\label{caratterizzazione-sequenziale}

Per \(x_0\in D'\) e \(L\in\overline{\mathbb R}\):

\[
\lim_{\substack{x\to x_0\\x\in D}}f(x)=L
\]

\[
\Longleftrightarrow
\forall(a_n)\subseteq D\setminus\{x_0\}:
a_n\to x_0
\Rightarrow
f(a_n)\to L.
\]

\[
\exists(a_n),(b_n)\subseteq D\setminus\{x_0\},
\quad
a_n,b_n\to x_0,
\]

\[
f(a_n)\to L_1,
\qquad
f(b_n)\to L_2,
\qquad
L_1\ne L_2
\Longrightarrow
\nexists\lim_{x\to x_0}f(x).
\]

\[
\delta_n=\dfrac{1}{n+1}\downarrow0.
\]

Approfondimento: \href{https://ingegnerismo.it/matematica/caratterizzazione-sequenziale-del-limite/}{caratterizzazione sequenziale del limite}.

Esercizi svolti: \href{https://ingegnerismo.it/matematica/caratterizzazione-sequenziale-del-limite-esercizi/}{caratterizzazione sequenziale del limite}.

\section{Limiti laterali}\label{limiti-laterali}

Le definizioni seguenti richiedono che \(x_0\) sia punto di accumulazione di \(D\cap(-\infty,x_0)\), rispettivamente di \(D\cap(x_0,+\infty)\).

\[
\lim_{\substack{x\to x_0^-\\x\in D}}f(x)=L
\Longleftrightarrow
\forall\varepsilon>0,\ \exists\delta>0,\ \forall x\in D:\
x_0-\delta<x<x_0
\Rightarrow
|f(x)-L|<\varepsilon,
\]

\[
\lim_{\substack{x\to x_0^+\\x\in D}}f(x)=L
\Longleftrightarrow
\forall\varepsilon>0,\ \exists\delta>0,\ \forall x\in D:\
x_0<x<x_0+\delta
\Rightarrow
|f(x)-L|<\varepsilon.
\]

Se \(x_0\) è punto di accumulazione di \(D\) da entrambi i lati:

\[
\lim_{\substack{x\to x_0\\x\in D}}f(x)=L
\Longleftrightarrow
\lim_{\substack{x\to x_0^-\\x\in D}}f(x)
=
\lim_{\substack{x\to x_0^+\\x\in D}}f(x)
=L.
\]

Approfondimento: \href{https://ingegnerismo.it/matematica/limiti-laterali/}{limiti laterali}.

\section{Limiti delle funzioni monotone}\label{limiti-delle-funzioni-monotone}

Sia \(f:I\to\mathbb R\) monotona e \(c\in\operatorname{int}I\).

\[
\begin{array}{c|c|c}
\text{monotonia}&f(c^-)&f(c^+)\\
\hline
\text{non decrescente}&\sup f(I\cap(-\infty,c))&\inf f(I\cap(c,+\infty))\\
\text{non crescente}&\inf f(I\cap(-\infty,c))&\sup f(I\cap(c,+\infty))
\end{array}
\]

\[
f\nearrow
\Longrightarrow
f(c^-)\le f(c)\le f(c^+),
\]

\[
f\searrow
\Longrightarrow
f(c^-)\ge f(c)\ge f(c^+).
\]

Se \(I\) è illimitato nelle direzioni indicate:

\[
\begin{array}{c|c|c}
\text{monotonia}&\lim_{x\to-\infty}f(x)&\lim_{x\to+\infty}f(x)\\
\hline
\text{non decrescente}&\inf f(I)&\sup f(I)\\
\text{non crescente}&\sup f(I)&\inf f(I)
\end{array}
\]

\[
f\text{ continua in }c
\Longleftrightarrow
f(c^-)=f(c)=f(c^+).
\]

Per \(f\) non decrescente:

\[
\lim_{x\to-\infty}f(x)=\inf f(I),
\qquad
\lim_{x\to+\infty}f(x)=\sup f(I),
\]

quando \(I\) è illimitato nella direzione corrispondente. Nei punti di discontinuità interna:

\[
J_f(c)=f(c^+)-f(c^-)>0.
\]

Esiste un'iniezione dell'insieme delle discontinuità in \(\mathbb Q\):

\[
c\longmapsto q_c,
\qquad
q_c\in\mathbb Q\cap\bigl(f(c^-),f(c^+)\bigr).
\]

Approfondimento: \href{https://ingegnerismo.it/matematica/limiti-delle-funzioni-monotone/}{limiti delle funzioni monotone}.

Esercizi svolti: \href{https://ingegnerismo.it/matematica/limiti-funzioni-monotone-esercizi/}{limiti delle funzioni monotone}.

\section{Algebra dei limiti e permanenza del segno}\label{algebra-dei-limiti-e-permanenza-del-segno}

Se \(f(x)\to\ell\), \(g(x)\to m\) con \(\ell,m\in\mathbb R\):

\[
f\pm g\to\ell\pm m,
\qquad
\lambda f\to\lambda\ell,
\]

\[
fg\to\ell m,
\qquad
|f|\to|\ell|,
\]

\[
\frac fg\to\frac\ell m,
\qquad m\ne0.
\]

\[
f(x)\to L>0
\Longrightarrow
f(x)>0\quad\text{definitivamente},
\]

\[
f(x)\to L<0
\Longrightarrow
f(x)<0\quad\text{definitivamente}.
\]

Come regola di limite, non come operazione aritmetica sui simboli infiniti:

\[
f(x)\to\ell\in\mathbb R,
\qquad
|g(x)|\to+\infty
\Longrightarrow
\frac{f(x)}{g(x)}\to0.
\]

Se \(f\to\ell\ne0\) e \(g\to0\):

\[
\begin{array}{c|cc}
&g\to0^+&g\to0^-\\
\hline
\ell>0&f/g\to+\infty&f/g\to-\infty\\
\ell<0&f/g\to-\infty&f/g\to+\infty
\end{array}
\]

Approfondimento: \href{https://ingegnerismo.it/matematica/algebra-dei-limiti/}{algebra dei limiti}, \href{https://ingegnerismo.it/matematica/permanenza-del-segno/}{permanenza del segno}.

\section{Composizione dei limiti}\label{composizione-dei-limiti}

Se

\[
f(x)\to\ell,
\qquad
g(y)\to m\quad(y\to\ell),
\]

e definitivamente

\[
f(x)\in D_g,
\qquad
f(x)\ne\ell,
\]

allora

\[
g(f(x))\to m.
\]

Se \(g\) è continua in \(\ell\):

\[
f(x)\to\ell
\Longrightarrow
g(f(x))\to g(\ell),
\]

senza richiedere \(f(x)\ne\ell\).

Approfondimento: \href{https://ingegnerismo.it/matematica/algebra-dei-limiti/}{algebra dei limiti}, \href{https://ingegnerismo.it/matematica/limite-di-funzione/}{limite di funzione}.

\section{Teorema dei carabinieri}\label{teorema-dei-carabinieri}

\[
f(x)\le g(x)\le h(x)\quad\text{definitivamente},
\]

\[
f(x)\to L,
\qquad
h(x)\to L
\Longrightarrow
g(x)\to L.
\]

Per \(L\in\mathbb R\), una forma equivalente particolarmente utile è:

\[
|g(x)-L|\le r(x),
\qquad
r(x)\to0
\Longrightarrow
g(x)\to L.
\]

Versioni infinite:

\[
f\le g,
\qquad
f\to+\infty
\Longrightarrow
g\to+\infty,
\]

\[
g\le h,
\qquad
h\to-\infty
\Longrightarrow
g\to-\infty.
\]

Approfondimento: \href{https://ingegnerismo.it/matematica/teorema-dei-carabinieri/}{teorema dei carabinieri}.

\section{Funzioni razionali all'infinito}\label{funzioni-razionali-allinfinito}

Siano

\[
P(x)=a_mx^m+\cdots,
\qquad
Q(x)=b_nx^n+\cdots,
\qquad a_m b_n\ne0.
\]

Per \(x\to\pm\infty\):

\[
\frac{P(x)}{Q(x)}
\sim
\frac{a_m}{b_n}x^{m-n}.
\]

In particolare, per \(x\to+\infty\):

\begin{longtable}[]{@{}ll@{}}
\toprule\noalign{}
Gradi & Limite di \(P/Q\) \\
\midrule\noalign{}
\endhead
\bottomrule\noalign{}
\endlastfoot
\(m<n\) & \(0\) \\
\(m=n\) & \(a_m/b_n\) \\
\(m>n\) & si studia il segno di \((a_m/b_n)x^{m-n}\) \\
\end{longtable}

Per \(x\to-\infty\) occorre considerare anche la parità di \(m-n\).

\section{Forme indeterminate}\label{forme-indeterminate}

Le scritture seguenti non sono risultati: indicano che le sole informazioni sui limiti dei fattori non bastano a determinare il limite dell'espressione.

\[
\frac00,
\qquad
\frac\infty\infty,
\qquad
0\cdot\infty,
\qquad
\infty-\infty,
\qquad
0^0,
\qquad
1^\infty,
\qquad
\infty^0.
\]

Trasformazioni standard:

\[
f\,g=\frac f{1/g}=\frac g{1/f},
\]

\[
f-g=\frac{f^2-g^2}{f+g}
\quad\text{oppure}\quad
f-g=\frac{1/g-1/f}{1/(fg)},
\]

nei rispettivi domini.

Per \(f(x)>0\):

\[
f(x)^{g(x)}
=\exp(g(x)\ln f(x)).
\]

Approfondimento: \href{https://ingegnerismo.it/matematica/forme-indeterminate/}{forme indeterminate}.

\section{Potenze con base ed esponente variabili}\label{potenze-con-base-ed-esponente-variabili}

Per una forma \(f(x)^{g(x)}\) con \(f(x)>0\) definitivamente, porre

\[
H(x)=g(x)\ln f(x).
\]

Se \(H(x)\to L\in\overline{\mathbb R}\), allora

\[
f(x)^{g(x)}=e^{H(x)}\to e^L,
\]

con le convenzioni \(e^{-\infty}=0\) ed \(e^{+\infty}=+\infty\).

Caso fondamentale:

\[
u(x)\to0,
\qquad
v(x)u(x)\to L
\Longrightarrow
[1+u(x)]^{v(x)}\to e^L,
\]

purché \(1+u(x)>0\) definitivamente.

\section{Limiti notevoli}\label{limiti-notevoli}

Per \(x\to0\):

\[
\frac{\sin x}{x}\to1,
\qquad
\frac{1-\cos x}{x^2}\to\frac12,
\qquad
\frac{\tan x}{x}\to1,
\]

\[
\frac{\arcsin x}{x}\to1,
\qquad
\frac{\arctan x}{x}\to1,
\]

\[
\frac{e^x-1}{x}\to1,
\qquad
\frac{a^x-1}{x}\to\ln a
\quad(a>0),
\]

\[
\frac{\ln(1+x)}{x}\to1,
\qquad
\frac{(1+x)^\alpha-1}{x}\to\alpha,
\]

\[
(1+x)^{1/x}\to e.
\]

Per \(x\to+\infty\):

\[
\left(1+\frac ax\right)^x\to e^a,
\]

\[
\frac{\ln x}{x^\alpha}\to0,
\qquad \alpha>0,
\]

\[
\frac{x^\alpha}{a^x}\to0,
\qquad a>1,
\qquad \alpha>0.
\]

Le formule trigonometriche richiedono gli angoli in radianti.

Approfondimento: \href{https://ingegnerismo.it/matematica/limiti-notevoli/}{limiti notevoli}.
